\documentclass{article}

\usepackage[utf8]{inputenc}
\usepackage[T1]{fontenc}
\usepackage{helvet}
\usepackage{geometry}
\usepackage{xcolor}
\usepackage{eso-pic}
\usepackage{fancyhdr}
\usepackage{parskip}
\usepackage{microtype}
\usepackage[english, ukrainian]{babel}
\usepackage{url}
\usepackage{amsmath}
\usepackage{amssymb}
\usepackage{amsthm}
\usepackage{longtable}
\usepackage{booktabs}
\usepackage{hyperref}
\usepackage{titlesec}

\definecolor{nato}{HTML}{293757}
\definecolor{footergray}{gray}{0.65}

\geometry{a4paper,left=3cm,right=3cm,top=3cm,bottom=3cm}
\renewcommand{\familydefault}{\sfdefault}
\setcounter{secnumdepth}{3}

\titleformat{\section}
  {\color{nato}\Huge\bfseries}{\thesection.}{1em}{}
\titlespacing*{\section}{0pt}{30pt}{12pt}

\titleformat{\subsection}
  {\color{nato}\LARGE\bfseries}{\thesubsection.}{1em}{}
\titlespacing*{\subsection}{0pt}{20pt}{8pt}

\titleformat{\subsubsection}
  {\color{nato}\Large\bfseries}{\thesubsubsection.}{1em}{}
\titlespacing*{\subsubsection}{0pt}{10pt}{6pt}

\fancyhf{}
\fancyhead[R]{\small\color{footergray} 23 квітня 2026}
\fancyfoot[L]{\small\color{footergray} Технічні вимоги ЄСІКС. VI. Суддівство}

\fancyfoot[R]{\small\color{footergray}\thepage}

\newcommand\HUGE[1]{\fontsize{40}{40}\selectfont #1}

\renewcommand{\headrulewidth}{0pt}
\renewcommand{\footrulewidth}{0.4pt}
\renewcommand{\footrule}{\color{footergray}\hrule width \textwidth height 0.4pt}

\begin{document}

\pagestyle{fancy}

\begin{titlepage}
  \AddToShipoutPictureBG*{\AtPageLowerLeft{\color{nato}\rule{\paperwidth}{\paperheight}}}
  \color{white}\thispagestyle{empty}\vspace*{4cm}\centering
  {\Huge  \textbf{Технічні вимоги ЄСІКС} \par} \vspace{2cm}
  {\HUGE  \textbf{VI. Суддівство} \par} \vspace{2.5cm}
  {\LARGE \textbf{VI-1. Облік суддів та суддівських посад} \par}
  {\LARGE \textbf{VI-2. Управління процедурами суддівської кар'єри} \par}
  {\LARGE \textbf{VI-3. Регулярне оцінювання суддів} \par}
  {\LARGE \textbf{VI-4. Суддівське досьє} \par}
  {\LARGE \textbf{VI-5. Верифікація} \par} \vspace{2.5cm}
  \vfill
  {\large 2026 \copyright\ Інформаційні Судові Системи \par}
\end{titlepage}

\newpage

\begin{center}
  {\Huge\color{nato}\bfseries Зміст\par}
\end{center}
\vspace{40pt}
\tableofcontents

\begin{abstract}
ТЕХНІЧНІ ВИМОГИ
НА СТВОРЕННЯ ЗАСОБУ ІНФОРМАТИЗАЦІЇ - КОМПОНЕНТІВ “СУДДІВСТВО”, ЖИТТЄВОГО ЦИКЛУ ПОСАДИ СУДДІ
ЄДИНОЇ СУДОВОЇ ІНФОРМАЦІЙНО-КОМУНІКАЦІЙНОЇ СИСТЕМИ.
\end{abstract}

\newpage
\section*{Вступна частина}

\subsection*{Умовні скорочення та визначення}

\begin{table}[ht]
\caption{Умовні скорочення та визначення}
\begin{tabular}{p{4cm}p{10cm}}
\hline
\textbf{Термін} & \textbf{Значення} \\
\hline
API & Автоматизовані програмні інтеграції (application programming interface) \\
БД & База Даних \\
БП & Бізнес Процес \\
КЕП & Кваліфікований електронний підпис (X.509) \\
ЄСІКС & Єдина судова інформаційно-комунікаційна система \\
ВККС & Вища кваліфікаційна комісія суддів України \\
ВРП & Вища рада правосуддя \\
НАЗК & Національне агентство з питань запобігання корупції \\
\hline
\end{tabular}
\end{table}

\subsection*{Загальні положення}

Технічні вимоги розроблено на основі «Концепції ЄСІКС» від 2024 року. Вимоги до суддівського досьє (розділ VI-4) сформульовано на рівні сутностей, інваріантів та процесів відповідно до політики документування. Нефункціональні вимоги визначаються централізованими сервісами (foundations.tex).

\newpage
\section{Облік суддів та суддівських посад}

Система забезпечує централізований облік судів, суддівських посад та фізичних осіб.

\subsection{Вимоги до сутностей}

\subsubsection{Суд та суддівська посада}
Сутність «Суддівська посада» (Judicial Position) містить реквізити посади відповідно до структури органів судової влади.

\begin{table}[ht]
\caption{Атрибути сутності «Суддівська посада»}
\begin{tabular}{p{5cm}p{9cm}}
\hline
\textbf{Атрибут} & \textbf{Опис} \\
\hline
ID посади & Унікальний системний ідентифікатор посади \\
Назва суду & Повне найменування судової установи \\
Рівень та спеціалізація & Ранг суду та його спеціалізація (касаційний, апеляційний тощо) \\
Статус вакантності & Поточний стан (вакантна, зайнята, тимчасово вакантна) \\
Адміністративна посада & Ознака та тип адмінпосади (голова суду, заступник) \\
Історія руху & Журнал призначень, переведень та звільнень за посадою \\
\hline
\end{tabular}
\end{table}

\subsubsection{Фізична особа (кандидат / суддя)}
Базова сутність «Фізична особа» (Physical Person) є спільною для кандидатів і суддів, забезпечуючи наскрізну ідентифікацію.

\begin{table}[ht]
\caption{Атрибути сутності «Фізична особа»}
\begin{tabular}{p{5cm}p{9cm}}
\hline
\textbf{Атрибут} & \textbf{Опис} \\
\hline
ПІБ & Прізвище, ім'я та по батькові особи \\
РНОКПП & Реєстраційний номер облікової картки платника податків \\
Обліковий статус & Поточна роль у системі (Кандидат, Суддя) \\
Статус перебування & Стан особи (діючий, звільнений, відсторонений, у відставці) \\
Контактні дані & Електронна пошта, номер телефону, адреса для листування \\
Посилання на досьє & Зв'язок з агрегатною сутністю «Досьє» \\
\hline
\end{tabular}
\end{table}

\subsection{Вимоги до процесів}

\begin{itemize}
\item Створення та ведення майстер-реєстрів судів і суддівських посад.
\item Реєстрація нового або діючого судді зі створенням персонального досьє.
\item Автоматичне створення запису про кандидата на підставі заяви на участь у процедурі.
\end{itemize}

\newpage
\section{Адміністрування процедур}

Система підтримує повний життєвий цикл процедур добору, конкурсу та кар’єрних переміщень суддів.

\subsection{Вимоги до сутностей}

\subsubsection{Процедура суддівської кар'єри}
Сутність «Процедура» (Procedure) — конкурс, добір або кваліфоцінювання з визначеним набором етапів та правил.

\begin{table}[ht]
\caption{Атрибути сутності «Процедура суддівської кар'єри»}
\begin{tabular}{p{5cm}p{9cm}}
\hline
\textbf{Атрибут} & \textbf{Опис} \\
\hline
ID процедури & Унікальний номер та шифр процедури \\
Тип процедури & Категорія (добір, конкурс, переведення, оцінювання) \\
Статус & Поточна фаза (проєкт, оголошена, триває, завершена) \\
Рішення Комісії & Реквізити нормативного акта про початок процедури \\
Етапи процедури & Перелік та послідовність кроків (іспит, співбесіда тощо) \\
Предмет & Кількість вакансій або конкретний перелік посад \\
\hline
\end{tabular}
\end{table}

\subsubsection{Учасник процедури}
Учасник процедури (Participant) — фізична особа, допущена до участі в конкретній процедурі з фіксацією статусу проходження.

\begin{table}[ht]
\caption{Атрибути сутності «Учасник процедури»}
\begin{tabular}{p{5cm}p{9cm}}
\hline
\textbf{Атрибут} & \textbf{Опис} \\
\hline
Суб'єкт & Посилання на профіль фізичної особи \\
Статус участі & Стан (заяву подано, допущено, відмовлено, переможець) \\
Результати етапів & Деталізація балів та оцінок за кожен пройдений етап \\
Рейтинговий бал & Сумарний показник для формування черговості \\
\hline
\end{tabular}
\end{table}

\subsection{Вимоги до процесів}

\begin{itemize}
\item Створення, редагування, оголошення та скасування процедур.
\item Подання учасником звернення щодо участі.
\item Проведення іспитів, декодування результатів.
\item Проведення засідань, формування порядку денного.
\item Формування рейтингів та визначення переможців.
\end{itemize}

\newpage
\section{Регулярне оцінювання суддів}

\subsection{Вимоги до сутностей}

\subsubsection{Оцінювання}
Сутність «Оцінювання» (Evaluation) пов’язана з суддею, періодом та результатами регулярного або кваліфікаційного оцінювання.

\begin{table}[ht]
\caption{Атрибути сутності «Оцінювання»}
\begin{tabular}{p{5cm}p{9cm}}
\hline
\textbf{Атрибут} & \textbf{Опис} \\
\hline
Об'єкт оцінювання & Посилання на досьє судді \\
Період & Часовий проміжок, за який проводиться оцінювання \\
Анкети & Набір заповнених форм (суддею, викладачами НШСУ тощо) \\
Результати & Підсумкові бали та висновки щодо відповідності посаді \\
Заперечення & Дані про подані заперечення та результати їх розгляду \\
\hline
\end{tabular}
\end{table}

\subsection{Вимоги до процесів}

\begin{itemize}
\item Визначення учасників та контроль умов участі.
\item Проведення оцінювання та узагальнення результатів.
\item Подання та розгляд заперечень судді.
\item Реєстрація та припинення участі громадського об’єднання.
\end{itemize}

\newpage
\section{Досьє судді (кандидата)}

\subsection{Загальна модель досьє}

Суддівське досьє (Dossier) реалізується як агрегатна сутність, що об'єднує структуровані дані та електронні документи.

\begin{table}[ht]
\caption{Атрибути сутності «Суддівське досьє»}
\begin{tabular}{p{5cm}p{9cm}}
\hline
\textbf{Атрибут} & \textbf{Опис} \\
\hline
Власник досьє & Посилання на фізичну особу \\
Тип досьє & Категорія (досьє кандидата, суддівське досьє) \\
Структура розділів & Набір логічних блоків згідно з нормативними вимогами \\
Публічна версія & Сформований набір даних для оприлюднення (знеособлений) \\
Журнал аудиту & Повна історія переглядів та змін у досьє \\
\hline
\end{tabular}
\end{table}

Інваріанти досьє:
\begin{itemize}
\item Єдине досьє протягом усього життєвого циклу (кандидат -> суддя).
\item Автоматичне створення при першому зверненні або реєстрації.
\item Досьє є джерелом істини для всіх пов’язаних процесів.
\end{itemize}

\subsubsection{Документи у складі досьє}
Кожен документ певного виду автоматично підключається до відповідного розділу досьє.

\subsection{Вимоги до процесів досьє}

\begin{itemize}
\item Перегляд власного досьє в особистому кабінеті.
\item Візуалізація структури та розділів.
\item Формування публічних версій досьє кандидатів.
\item Надання тимчасового доступу уповноваженим користувачам.
\end{itemize}

\newpage
\section{Верифікація}

Розділ охоплює функціональність модулів «Аналітика даних» (5.10) та підетапів перевірок (5.3.4, 5.3.5) з оригінального ТВ.

\subsection{Вимоги до сутностей}

\subsubsection{Перевірка}
Сутність «Перевірка» (Verification) — формалізований запис результатів контрольних заходів щодо судді або кандидата.

\begin{table}[ht]
\caption{Атрибути сутності «Перевірка»}
\begin{tabular}{p{5cm}p{9cm}}
\hline
\textbf{Атрибут} & \textbf{Опис} \\
\hline
Тип перевірки & Вид (доброчесність, родинні зв'язки, спеціальна тощо) \\
Суб'єкт & Посилання на досьє особи \\
Статус & Поточний стан (ініційовано, дані отримано, завершено) \\
Результати & Агреговані дані зіставлення (структурований набір даних) \\
Документ-висновок & Посилання на документ у «Документ-сервісі» з КЕП \\
\hline
\end{tabular}
\end{table}

\subsubsection{Резерв кандидатів}
Сутність «Резерв кандидатів» (Candidate Reserve) забезпечує облік осіб за результатами рейтингів.

\begin{table}[ht]
\caption{Атрибути сутності «Резерв кандидатів»}
\begin{tabular}{p{5cm}p{9cm}}
\hline
\textbf{Атрибут} & \textbf{Опис} \\
\hline
Кандидат & Посилання на досьє фізичної особи \\
Дата формування & Дата включення до резерву \\
Строк перебування & Граничний термін (параметр системи) \\
Підстава зміни & Посилання на рішення Комісії \\
\hline
\end{tabular}
\end{table}

\newpage
\subsection{Вимоги до процесів}

\subsubsection{Ризик-орієнтовані перевірки (5.10.1)}
\begin{enumerate}
    \item Формування запиту до Модуля збору інформації ЄСІКС.
    \item Отримання та аналітична обробка даних на боці підсистеми.
    \item Формування документа-висновку за шаблоном «Документ-сервісу».
\end{enumerate}

\subsubsection{Спеціальна перевірка (5.3.5)}
\begin{enumerate}
    \item Ініціювання процедури уповноваженою особою ВККС.
    \item Формування пакету запитів (підписання КЕП однією дією).
    \item Подання документів особою та фіксація результатів перевірки.
\end{enumerate}

\subsubsection{Управління резервом (5.11.1)}
\begin{enumerate}
    \item Автоматичне додавання учасників до резерву на основі рейтингів.
    \item Вилучення або продовження за рішенням Комісії.
    \item Експорт даних резерву (PDF, XLSX, CSV).
\end{enumerate}

\newpage
\section*{Висновок}

Розділ VI «Суддівство» забезпечує єдиний цифровий контур життєвого циклу посади судді. Суддівське досьє (VI-4) визначено як формальна агрегатна сутність з чіткими інваріантами. Усі процеси інтегруються через базові сервіси ЄСІКС (foundations.tex).

\begin{thebibliography}{9}
\bibitem{esiks} Концепція ЄСІКС. \url{https://court.gov.ua/storage/portal/dsa/news/Kontseptsia%20ESIKS.pdf}. 2024
\bibitem{policy} Політики ЄСІКС. \texttt{policy.pdf}. 2026
\bibitem{iso29148} ISO/IEC/IEEE 29148:2018. Systems and software engineering — Requirements engineering.
\bibitem{iso19510} ISO/IEC 19510:2013. Business Process Model and Notation (BPMN).
\bibitem{dstu:42010} ДСТУ 42010:2018. \newblock Інженерія систем і програмних засобів. Опис архітектури.
\bibitem{law:personal-data} Закон України «Про захист персональних даних».
\bibitem{esiks} Концепція ЄСІКС. 2024.
\end{thebibliography}

\end{document}
